\section{Introduction}

Large language models (LLMs) have demonstrated remarkable potential for clinical reasoning, achieving expert-level performance on medical licensing examinations~\cite{singhal2023large} and showing strong capabilities across a diverse range of medical tasks~\cite{bedi2025medhelm, wu2025towards, wang2025capabilities}. Recent works have extended these capabilities to agentic settings, where LLMs interact with electronic health record (EHR) databases to answer clinically meaningful questions~\cite{liao2026agentehr, tang2024medagents}. These agents typically generate code or queries to extract structured information from EHR systems and reason over retrieved data to support clinical decisions.

A common approach in these systems is to equip agents with a user-designed toolbox: a curated set of domain-specific functions that the agent can invoke to handle recurring clinical computations. While this strategy improves reliability over purely zero-shot approaches, it does not scale. Designing, validating, and maintaining such toolboxes requires sustained effort from clinical and engineering experts, and the scope of the toolbox is ultimately limited by what its designers anticipated. Moreover, when agents re-execute lengthy tool chains for every complex query, the token cost compounds quickly. This is a practical concern in high-stakes medical deployments where latency and cost matter.

Beyond the toolbox scalability problem, there is a more fundamental issue. Many clinically meaningful questions are simply not well-defined under a zero-shot evaluation regime. Clinical concepts such as sepsis determination, organ failure scoring, and vasopressor thresholds are not universal; their precise operationalizations vary across institutions and evolve with clinical guidelines. An agent answering such questions without access to the relevant institutional definitions faces an underspecified problem, and any benchmark grading its output against a single fixed ground truth will produce misleading accuracy estimates.

These two problems (the fragility of static toolboxes and the ambiguity of zero-shot clinical definitions) point toward a different evaluation paradigm. Rather than asking \emph{how well does an agent answer clinical questions given a fixed toolbox?}, a more scalable and meaningful question is: \emph{given limited verification data, how well can LLMs generate modular, composable libraries for processing EHR data?} If LLMs can reliably synthesize such libraries from clinical guidelines, the bottleneck of expert-curated toolboxes is removed, and the generated code itself encodes the institutional definitions needed to make questions well-posed.

We introduce \textbf{\ours{}}, a benchmark designed to evaluate exactly this capability. \ours{} is grounded in MIMIC-IV~\cite{johnson2020mimic} and spans 63 clinical concepts across nine domains, including severity scores, sepsis criteria, organ failure indices, and ICU treatments. It contains 259 question variants and 63{,}000 QA pairs with difficulty stratified by \emph{compositional dependency depth}, defined as the number of concept-to-concept dependencies an agent must resolve. Crucially, \ours{} also includes a \emph{longitudinal reasoning} component that evaluates agents on questions tracking patient state across the full course of ICU admissions. \junjie{Highlight the key data statistics on the longitudinal agentic task. Should we call it ``longitudinal agentic automation'' instead of ``logitudinal reasoning''?} Both of these properties make the questions genuinely difficult to resolve zero-shot: compositional questions require chaining multi-step computations that naive approaches repeatedly reinvent, while longitudinal questions require consistent state-tracking logic that zero-shot code generation rarely produces reliably. Together, they establish a setting where the quality of an agent's underlying function library, rather than its in-context reasoning alone, is the primary determinant of performance.

To establish a strong baseline, we propose an offline autoformalization pipeline that converts natural-language clinical guidelines into a reusable, QA-verified Python function library via iterative LLM generation. This library is constructed once using limited verification labels, then reused across all queries, eliminating redundant computation and enforcing consistent definitions across runs.

Our contributions are as follows:
\begin{itemize}
\item \textbf{\ours\ (Benchmark)}: 63 clinical concepts, 259 question variants, and 63{,}000 QA pairs derived from MIMIC-IV, with difficulty stratified by compositional dependency depth and a longitudinal reasoning component. The benchmark is designed to surface the failure modes of static toolboxes and zero-shot code generation on questions that are difficult to define without institutional context.
\item \textbf{Benchmark Analysis (Findings)}: A systematic characterization of where and why current LLMs fail on compositional and longitudinal clinical reasoning, including analyses by difficulty level, question type, model scale, and EHR schema familiarity.
\item \textbf{Clinical Autoformalization (Baseline)}: An automated pipeline that transforms natural-language clinical guidelines into a verified Python function library using limited verification data, demonstrating a scalable alternative to user-designed toolboxes and serving as a strong baseline for future work on \ours{}.
\end{itemize}
